\documentclass[8pt]{extarticle}

\usepackage{diplomstyle}

\geometry{
	a4paper,
	paperheight=140mm,
	paperwidth=140mm,
	left=0.5cm,
	right=0.5cm,
	top=0.5cm,
	bottom=0.5cm,
	footskip=0.8cm
}

\usepackage{booktabs}
\usepackage{array}
\usepackage{makecell}
\usepackage{amsmath}
\usepackage{amssymb}
\usepackage{ragged2e}
\usepackage{longtable}
\usepackage{esint}
\usepackage{draftwatermark}
\usepackage{tikz}
\usetikzlibrary{patterns,arrows.meta,shapes.geometric}
\usepackage{circuitikz}
\usepackage{cancel}

\DeclareMathOperator{\re}{\operatorname{Re}}
\DeclareMathOperator{\im}{\operatorname{Im}}
\DeclareMathOperator{\Epsilon}{\mathcal{E}}
\DeclareMathOperator{\Div}{\operatorname{div}}
\DeclareMathOperator{\Rot}{\operatorname{rot}}
\DeclareMathOperator{\Grad}{\operatorname{grad}}
\DeclareMathOperator{\R}{\mathbb{R}}

\makeatletter
\newcommand\incircbin{\mathpalette\@incircbin}
\newcommand\@incircbin[2]{%
	\mathbin{\ooalign{\hidewidth$#1#2$\hidewidth\crcr$#1\bigcirc$}}%
}
\newcommand{\oeq}{\incircbin{=}}
\makeatother

\newcommand{\pp}[2]{\frac{\partial #1}{\partial #2}}
\newcommand{\dd}[2]{\frac{d {#1}}{d {#2}}}
\newcommand{\ppp}[2]{\frac{\partial^2 #1}{\partial #2^2}}
\newcommand{\ddd}[2]{$\frac{d^2 {#1}}{d {#2}^2}$}

\renewcommand{\theadfont}{\normalsize}
\newcolumntype{L}[1]{>{\raggedright\arraybackslash}p{#1}}
\newcolumntype{C}[1]{>{\centering\arraybackslash}p{#1}}

\title{Ответы на вопросы \\ <<Теория информации и кодирования>>}
\author{Скороходов С.А., 438Б группа}
\date{ }

\SetWatermarkText{\textbf{SerKin0}}
\SetWatermarkColor[gray]{0.95}
\SetWatermarkScale{1}

\def\tquastion#1{\newpage\addcontentsline{toc}{subsection}{#1}\begin{librarybox}#1 \end{librarybox}}

\begin{document}
	\setlength{\abovedisplayskip}{3px}
	\setlength{\belowdisplayskip}{3px}
	\setlength{\abovedisplayshortskip}{3px}
	\setlength{\belowdisplayshortskip}{3px}
	\maketitle
	Ответы на вопросы по дисциплине <<Теория информации и кодирования>> подготовлены для коллоквиума 2026--2027 учебного года для студентов радиофизического факультета. Данный материал предназначен для направления подготовки <<Информационная безопасность телекоммуникационных систем>> и разработан на основе лекций и практических занятий, проводимых преподавателем Мальцев Александр Александрович.
	
	Благодарность в помощи составления ответов выражается Помелову Сергею Дмитриевичу.
	\tableofcontents
	
	\tquastion{1. Интегральная и дифференциальная функция распределения случайной величины и их свойства.}
	
	Случайная величина $X$ --- числовая функция исхода случайного опыта.
	
	\textbf{Интегральная функция распределения:}
	\[
	\boxed{F_X(x)=P\{X\leqslant x\}.}
	\]
	Показывает вероятность того, что случайная величина примет значение, не превосходящее $x$. Свойства: $0\leqslant F_X(x)\leqslant1$; неубывает; непрерывна справа; $F_X(-\infty)=0$, $F_X(+\infty)=1$.
	\[
	\boxed{P\{a<X\leqslant b\}=F_X(b)-F_X(a).}
	\]
	
	\textbf{Дифференциальная функция распределения (плотность вероятности)} существует для абсолютно непрерывной случайной величины:
	\[
	\boxed{
		W_X(x)=\frac{dF_X(x)}{dx},\qquad
		F_X(x)=\int_{-\infty}^{x}W_X(x')\,dx'.
	}
	\]
	Физический смысл: $W_X(x)\,dx$ --- вероятность попадания случайной величины в малый интервал $[x,x+dx]$.
	
	Свойства:
	\[
	\boxed{
		W_X(x)\geqslant0,\qquad
		\int_{-\infty}^{+\infty}W_X(x)\,dx=1.
	}
	\]
	\[
	P\{a<X<b\}=\int_a^b W_X(x)\,dx.
	\]
	Плотность не является вероятностью в отдельной точке; для непрерывного распределения $P\{X=x\}=0$. В дискретном случае вместо обычной плотности используются вероятностные массы (или $\delta$-функции).
	
	\tquastion{2. Понятие стационарности случайного процесса: (а) в узком смысле; (б) в широком смысле.}
	
	\textbf{(а) Стационарность в узком смысле (строгая).}
	
	Все $n$-мерные плотности вероятности инвариантны относительно произвольного сдвига времени:
	\[
	\boxed{
		W_n(x_1,t_1;\ldots;x_n,t_n)
		=
		W_n(x_1,t_1+\tau;\ldots;x_n,t_n+\tau).
	}
	\]
	Следствия: $W(x,t)=W(x)$; двумерная плотность зависит только от $\tau=t_2-t_1$; среднее и дисперсия постоянны; корреляционная и ковариационная функции зависят только от $\tau$:
	\[
	K_x[t_1,t_2]=K_x(\tau),\qquad B_x[t_1,t_2]=B_x(\tau).
	\]
	
	Для действительного стационарного процесса корреляционная и ковариационная функции чётны:
	\[
	K_x(\tau)=K_x(-\tau),\qquad B_x(\tau)=B_x(-\tau).
	\]
	
	\textbf{(б) Стационарность в широком смысле.}
	
	Достаточно выполнения двух условий:
	\[
	\boxed{
		\langle x(t)\rangle=m_x=\mathrm{const},\qquad
		K_x[t_1,t_2]=K_x(t_2-t_1).
	}
	\]
	Эквивалентно, при постоянном среднем, ковариационная функция зависит только от $\tau$.
	
	Строгая стационарность влечёт стационарность в широком смысле, обратное в общем случае неверно. Для гауссовских процессов оба определения эквивалентны.
	
	Пример: $x(t)=A_0\cos(\omega_0t+\varphi)$ при равномерной случайной фазе $\varphi\in[-\pi,\pi]$ является стационарным процессом:
	\[
	\langle x(t)\rangle=0,\qquad
	\boxed{K_x(\tau)=\frac{A_0^2}{2}\cos(\omega_0\tau).}
	\]
	
	\tquastion{3. Выражения для n-моментной плотности вероятностей случайного процесса через условные плотности вероятностей: общий случай; марковский процесс; совершенно случайный процесс.}
	
	\textbf{Общий случай} (цепное правило):
	\[
	\boxed{
		\begin{aligned}
			W_n={}&W(x_1,t_1)W(x_2,t_2|x_1,t_1)\\
			&\cdot W(x_3,t_3|x_1,t_1;x_2,t_2)\cdots\\
			&\cdot W(x_n,t_n|x_1,t_1;\ldots;x_{n-1},t_{n-1}).
		\end{aligned}
	}
	\]
	
	\textbf{Марковский процесс} (процесс без последействия): будущее состояние зависит только от последнего известного значения:
	\[
	W(x_n,t_n|x_1,t_1;\ldots;x_{n-1},t_{n-1})
	=
	W(x_n,t_n|x_{n-1},t_{n-1}).
	\]
	Тогда
	\[
	\boxed{
		W_n=W(x_1,t_1)\prod_{i=2}^{n}
		W(x_i,t_i|x_{i-1},t_{i-1}).
	}
	\]
	
	\textbf{Совершенно случайный процесс:} значения процесса в любые несовпадающие моменты времени статистически независимы:
	\[
	\boxed{
		W_n(x_1,t_1;\ldots;x_n,t_n)
		=\prod_{i=1}^{n}W(x_i,t_i).
	}
	\]
	Для него условная плотность равна безусловной:
	$W(x_n,t_n|x_1,t_1;\ldots;x_{n-1},t_{n-1})=W(x_n,t_n)$.
	
	Совершенно случайный процесс является частным случаем марковского.
	
	\tquastion{4. Поведение условной плотности вероятности $W(x_2,t_2|x_1,t_1)$ эргодического случайного процесса $x(t)$ при $\tau \to \infty$, $\tau \to 0$ ($\tau = t_2 - t_1$).}
	
	При $\tau\to0$ для процесса, непрерывного по вероятности, значение практически не успевает измениться:
	\[
	\boxed{
		\lim_{\tau\to0}W(x_2,t_2|x_1,t_1)=\delta(x_2-x_1).
	}
	\]
	
	При $\tau\to\infty$ для процесса с конечной памятью (обладающего свойством перемешивания) зависимость от начального значения исчезает:
	\[
	\boxed{
		\lim_{\tau\to\infty}W(x_2,t_2|x_1,t_1)=W(x_2,t_2).
	}
	\]
	Для стационарного процесса $W(x_2,t_2)=W(x_2)$.
	
	Физический смысл: при малом интервале процесс сохраняет начальное состояние, при большом --- забывает его. \textbf{Важно:} одна лишь эргодичность не гарантирует второй предел; требуется дополнительное условие затухания статистической зависимости (перемешивание).
	
	\tquastion{5. Определение квазидетерминированного случайного процесса, выражение для его одномерной плотности вероятностей.}
	
	Квазидетерминированный случайный процесс --- процесс, реализации которого описываются функцией заданного вида, содержащей один или несколько случайных параметров:
	\[
	x(t)=s(t,\lambda),
	\]
	где $\lambda$ --- случайная величина с известной плотностью $W_\lambda(\lambda)$.
	
	Одномерная плотность вероятности:
	\[
	\boxed{
		W(x,t)=\int_{-\infty}^{\infty}
		W_\lambda(\lambda)\delta[x-s(t,\lambda)]\,d\lambda.
	}
	\]
	Для $n$-мерной плотности:
	\[
	W_n=\int_{-\infty}^{\infty}W_\lambda(\lambda)
	\prod_{i=1}^{n}\delta[x_i-s(t_i,\lambda)]\,d\lambda.
	\]
	
	\textbf{Пример:} $x(t)=A_0\cos(\omega_0t+\varphi)$, где $\varphi$ равномерно распределена в $[-\pi,\pi]$:
	\[
	W_\varphi(\varphi)=
	\begin{cases}
		\frac{1}{2\pi},&|\varphi|\leqslant\pi,\\
		0,&|\varphi|>\pi.
	\end{cases}
	\]
	Подстановка в формулу с $\delta$-функцией даёт:
	\[
	\boxed{
		W(x,t)=
		\begin{cases}
			\dfrac{1}{\pi\sqrt{A_0^2-x^2}},&|x|\leqslant A_0,\\[5pt]
			0,&|x|>A_0.
		\end{cases}
	}
	\]
	Плотность не зависит от времени и частоты, имеет U-образный вид: около экстремумов $\pm A_0$ значения встречаются чаще, чем около нуля.
	
	\tquastion{6. Определение n-мерной характеристической функции случайного процесса и её связь с n-мерной плотностью вероятностей случайного процесса.}
	
	$n$-мерная характеристическая функция --- $n$-кратное преобразование Фурье от $n$-мерной плотности вероятности:
	\[
	\boxed{
		\Theta_x(u_1,t_1;\ldots;u_n,t_n)
		=
		\left\langle
		\exp\left(j\sum_{k=1}^{n}u_kx(t_k)\right)
		\right\rangle.
	}
	\]
	
	В интегральной форме:
	\[
	\Theta_x=
	\int_{\mathbb R^n}
	e^{j\sum_{k=1}^{n}u_kx_k}
	W_n(x_1,t_1;\ldots;x_n,t_n)\,d^nx.
	\]
	
	Обратное преобразование Фурье:
	\[
	\boxed{
		W_n=\frac{1}{(2\pi)^n}\int_{\mathbb R^n}
		\Theta_x(u_1,t_1;\ldots;u_n,t_n)
		e^{-j\sum_{k=1}^{n}u_kx_k}\,d^nu.
	}
	\]
	
	Основные свойства:
	\[
	\boxed{
		\Theta_x(0,\ldots,0)=1,\qquad |\Theta_x|\leqslant1.
	}
	\]
	\[
	\Theta_x(-u_1,\ldots,-u_n)=\Theta_x^*(u_1,\ldots,u_n).
	\]
	Свойство согласованности: для $m<n$
	\[
	\Theta_x(u_1,t_1;\ldots;u_m,t_m)
	=
	\Theta_x(u_1,t_1;\ldots;u_m,t_m;0,\ldots,0).
	\]
	Характеристическая функция однозначно определяет плотность вероятности.
	
	\tquastion{7. Определение среднего значения случайного процесса $\langle x(t)\rangle$, среднего значения функции случайного процесса $\langle f[x(t)]\rangle$, дисперсии случайного процесса.}
	
	Среднее значение (математическое ожидание) процесса:
	\[
	\boxed{
		\langle x(t)\rangle=
		\int_{-\infty}^{\infty}xW(x,t)\,dx.
	}
	\]
	
	Среднее значение произвольной функции процесса:
	\[
	\boxed{
		\langle f[x(t)]\rangle=
		\int_{-\infty}^{\infty}f(x)W(x,t)\,dx.
	}
	\]
	
	Дисперсия --- средний квадрат отклонения от среднего:
	\[
	\boxed{
		\sigma_x^2(t)=
		\left\langle(x(t)-\langle x(t)\rangle)^2\right\rangle
		=\langle x^2(t)\rangle-\langle x(t)\rangle^2.
	}
	\]
	Усреднение здесь выполняется по статистическому ансамблю, а не по времени.
	
	\tquastion{8. Моментные функции случайного процесса (определение).}
	
	Моментная функция $s$-го порядка --- среднее произведение значений процесса в $s$ моментах времени:
	\[
	\boxed{
		\alpha_s(t_1,\ldots,t_s)
		=\langle x(t_1)x(t_2)\cdots x(t_s)\rangle.
	}
	\]
	
	Через $s$-мерную плотность вероятности:
	\[
	\alpha_s=
	\int_{\mathbb R^s}x_1x_2\cdots x_s
	W_s(x_1,t_1;\ldots;x_s,t_s)\,d^sx.
	\]
	
	При совпадении всех моментов времени:
	$\alpha_s(t,\ldots,t)=\langle x^s(t)\rangle$.
	
	Связь с характеристической функцией:
	\[
	\boxed{
		\left.
		\frac{\partial^s\Theta_x}
		{\partial u_{k_1}\cdots\partial u_{k_s}}
		\right|_{\mathbf u=0}
		=j^s\alpha_s(t_{k_1},\ldots,t_{k_s}).
	}
	\]
	Первые два момента:
	\[
	\alpha_1(t)=m_x(t),\qquad
	\alpha_2(t_1,t_2)=K_x(t_1,t_2).
	\]
	
	\tquastion{9. Кумулянтные функции случайного процесса (определение).}
	
	Кумулянтные функции --- коэффициенты разложения логарифма характеристической функции в ряд Маклорена:
	\[
	\boxed{
		\ln\Theta_x=
		\sum_{s=1}^{\infty}\frac{j^s}{s!}
		\sum_{k_1,\ldots,k_s=1}^{n}
		\kappa_s(t_{k_1},\ldots,t_{k_s})
		u_{k_1}\cdots u_{k_s}.
	}
	\]
	
	Или через производные:
	\[
	\boxed{
		\kappa_s=
		\left.
		\frac{1}{j^s}
		\frac{\partial^s\ln\Theta_x}
		{\partial u_{k_1}\cdots\partial u_{k_s}}
		\right|_{\mathbf u=0}.
	}
	\]
	
	Связь с моментами:
	\[
	\boxed{
		\kappa_1(t)=m_x(t),\qquad
		\kappa_2(t_1,t_2)=
		\alpha_2(t_1,t_2)-\alpha_1(t_1)\alpha_1(t_2)=B_x(t_1,t_2).
	}
	\]
	Для одномерного случая четвёртый кумулянт:
	$\kappa_4=\mu_4-3\sigma_x^4$, где $\mu_4$ --- четвёртый центральный момент.
	
	Кумулянты независимых случайных величин складываются. Для гауссовского процесса все кумулянты порядка выше второго равны нулю.
	
	\tquastion{10. Корреляционная и ковариационная функции случайного процесса (определение).}
	
	\textbf{Корреляционная функция (нецентрированный момент второго порядка):}
	\[
	\boxed{
		K_x(t_1,t_2)=\langle x(t_1)x(t_2)\rangle.
	}
	\]
	
	\textbf{Ковариационная функция (центрированная корреляционная функция):}
	\[
	\boxed{
		\begin{aligned}
			B_x(t_1,t_2)
			&=\left\langle
			(x(t_1)-\langle x(t_1)\rangle)
			(x(t_2)-\langle x(t_2)\rangle)
			\right\rangle\\
			&=K_x(t_1,t_2)-m_x(t_1)m_x(t_2).
		\end{aligned}
	}
	\]
	
	При совпадающих аргументах:
	\[
	\boxed{B_x(t,t)=\sigma_x^2(t).}
	\]
	Для действительного процесса функции симметричны по аргументам:
	$K_x(t_1,t_2)=K_x(t_2,t_1)$, $B_x(t_1,t_2)=B_x(t_2,t_1)$.
	
	Неравенство Коши--Буняковского:
	\[
	\boxed{|B_x(t_1,t_2)|\leqslant\sigma_x(t_1)\sigma_x(t_2).}
	\]
	При стационарности обе функции зависят только от $\tau=t_2-t_1$.
	
	\tquastion{11. Определение коэффициента автокорреляции случайного процесса.}
	
	Коэффициент автокорреляции --- нормированная ковариационная функция:
	\[
	\boxed{
		R_x(t_1,t_2)=
		\frac{B_x(t_1,t_2)}
		{\sqrt{B_x(t_1,t_1)B_x(t_2,t_2)}}
		=\frac{B_x(t_1,t_2)}
		{\sigma_x(t_1)\sigma_x(t_2)}.
	}
	\]
	
	Для стационарного процесса:
	\[
	\boxed{
		R_x(\tau)=\frac{B_x(\tau)}{B_x(0)}
		=\frac{B_x(\tau)}{\sigma_x^2}.
	}
	\]
	
	Свойства: $R_x(0)=1$, $|R_x(\tau)|\leqslant1$; для действительного стационарного процесса $R_x(-\tau)=R_x(\tau)$. Коэффициент характеризует степень линейной связи между значениями процесса. Нулевой коэффициент означает некоррелированность, но не обязательно статистическую независимость.
	
	\tquastion{12. Определение коэффициента взаимной корреляции двух случайных процессов.}
	
	Для процессов $x(t)$ и $y(t)$ взаимная корреляционная функция:
	\[
	\boxed{K_{xy}(t_1,t_2)=\langle x(t_1)y(t_2)\rangle.}
	\]
	
	Взаимная ковариационная функция:
	\[
	\boxed{
		B_{xy}(t_1,t_2)=
		K_{xy}(t_1,t_2)-m_x(t_1)m_y(t_2).
	}
	\]
	
	Коэффициент взаимной корреляции:
	\[
	\boxed{
		R_{xy}(t_1,t_2)=
		\frac{B_{xy}(t_1,t_2)}
		{\sigma_x(t_1)\sigma_y(t_2)},\qquad |R_{xy}|\leqslant1.
	}
	\]
	
	Свойства симметрии:
	\[
	K_{xy}(t_1,t_2)=K_{yx}(t_2,t_1),\qquad
	B_{xy}(t_1,t_2)=B_{yx}(t_2,t_1).
	\]
	
	Для совместно стационарных процессов взаимные функции зависят только от $\tau$; при этом $R_{xy}(\tau)=R_{yx}(-\tau)$, но вообще не обязана выполняться чётность $R_{xy}(\tau)$.
	
	Статистическая независимость процессов влечёт их некоррелированность, обратное в общем случае неверно. Для совместно гауссовских процессов некоррелированность эквивалентна независимости.
	
	\tquastion{13. Отыскание среднего значения $\langle x(t)\rangle$ и дисперсии $\sigma_x^2$ случайного процесса по его корреляционной функции (для физических процессов с конечной памятью).}
	
	Для физического процесса с конечной памятью ковариационная функция стремится к нулю:
	\[
	\lim_{\tau\to\infty}B_x(\tau)=0.
	\]
	
	Учитывая, что $K_x(\tau)=B_x(\tau)+m_x^2$, получаем:
	\[
	\boxed{
		\lim_{\tau\to\infty}K_x(\tau)=m_x^2.
	}
	\]
	
	Следовательно:
	\[
	\boxed{
		m_x^2=\lim_{\tau\to\infty}K_x(\tau),\qquad
		\sigma_x^2=K_x(0)-\lim_{\tau\to\infty}K_x(\tau).
	}
	\]
	Знак среднего по одной корреляционной функции определить нельзя:
	$m_x=\pm\sqrt{\lim_{\tau\to\infty}K_x(\tau)}$.
	
	\tquastion{14. Эргодичность случайного процесса по отношению к среднему (необходимое и достаточное условие, достаточные условия).}
	
	Процесс эргодичен по отношению к среднему, если среднее по времени одной реализации при $T\to\infty$ стремится к среднему по статистическому ансамблю:
	\[
	\boxed{
		\lim_{T\to\infty}\frac{1}{T}\int_0^T x(t)\,dt
		=\langle x\rangle=m_x.
	}
	\]
	
	Рассмотрим временное среднее:
	\[
	\overline{x}_T=\frac{1}{T}\int_0^T x(t)\,dt,
	\qquad
	\langle\overline{x}_T\rangle=m_x.
	\]
	Его дисперсия для стационарного процесса:
	\[
	\begin{aligned}
		\sigma_{\overline{x}_T}^2
		&=\frac{1}{T^2}\int_0^T\int_0^T
		B_x(t_2-t_1)\,dt_1dt_2\\
		&=\frac{1}{T}\int_{-T}^{T}
		\left(1-\frac{|\tau|}{T}\right)B_x(\tau)\,d\tau\\
		&=\frac{2}{T}\int_0^T
		\left(1-\frac{\tau}{T}\right)B_x(\tau)\,d\tau.
	\end{aligned}
	\]
	
	\textbf{Необходимое и достаточное условие} среднеквадратической эргодичности:
	\[
	\boxed{
		\lim_{T\to\infty}
		\frac{1}{T}\int_{-T}^{T}
		\left(1-\frac{|\tau|}{T}\right)
		B_x(\tau)\,d\tau=0.
	}
	\]
	
	\textbf{Достаточное условие Слуцкого:}
	\[
	\boxed{
		\lim_{T\to\infty}\frac{1}{T}
		\int_0^T B_x(\tau)\,d\tau=0.
	}
	\]
	Более сильное, но простое для проверки условие:
	\[
	\boxed{\lim_{\tau\to\infty}B_x(\tau)=0.}
	\]
	Также достаточным является абсолютная интегрируемость:
	$\int_0^\infty|B_x(\tau)|\,d\tau<\infty$.
	
	\tquastion{15. Плотность вероятностей случайного процесса при безынерционном нелинейном преобразовании $y = f(x)$ для однозначной и многозначной обратной функции.}
	
	Безынерционное преобразование означает, что $y(t)=f[x(t)]$.
	
	Общая формула через $\delta$-функцию:
	\[
	\boxed{
		W_y(y,t)=\int_{-\infty}^{\infty}
		W_x(x,t)\delta[y-f(x)]\,dx.
	}
	\]
	
	\textbf{Однозначная обратная функция} $x=\varphi(y)$:
	\[
	\boxed{
		W_y(y,t)=W_x(\varphi(y),t)
		\left|\frac{d\varphi(y)}{dy}\right|
		=\frac{W_x(\varphi(y),t)}{|f'(\varphi(y))|}.
	}
	\]
	
	\textbf{Многозначная обратная функция:} если уравнение $f(x)=y$ имеет несколько простых корней $x_i=\varphi_i(y)$, то суммируем вклады всех ветвей:
	\[
	\boxed{
		W_y(y,t)=\sum_i W_x(\varphi_i(y),t)
		\left|\frac{d\varphi_i(y)}{dy}\right|
		=\sum_i\frac{W_x(\varphi_i(y),t)}
		{|f'(\varphi_i(y))|}.
	}
	\]
	
	Формулы следуют из сохранения вероятности. Они применимы в точках, где $f'(x_i)\neq0$; особые точки рассматривают по пределу.
	
	\tquastion{16. Гауссовский случайный процесс (определение). Вид n-мерной плотности вероятностей и характеристической функции.}
	
	Гауссовский случайный процесс --- процесс, для которого любая конечная совокупность отсчётов имеет совместное многомерное нормальное распределение. Эквивалентно, все кумулянты порядка $s\geqslant3$ равны нулю.
	
	Пусть $\mathbf m$ --- вектор математических ожиданий, $\mathbf B$ --- ковариационная матрица с элементами $B_{kl}=B_x(t_k,t_l)$. Тогда характеристическая функция:
	\[
	\boxed{
		\Theta_x(\mathbf u)=
		\exp\left(
		j\mathbf m^T\mathbf u-\frac12\mathbf u^T\mathbf B\mathbf u
		\right).
	}
	\]
	
	$n$-мерная плотность вероятностей при невырожденной ковариационной матрице:
	\[
	\boxed{
		\begin{aligned}
			W_n(\mathbf x)={}&
			\frac{1}{(2\pi)^{n/2}\sqrt{\det\mathbf B}}\\
			&\cdot\exp\left[
			-\frac12(\mathbf x-\mathbf m)^T
			\mathbf B^{-1}(\mathbf x-\mathbf m)
			\right].
		\end{aligned}
	}
	\]
	
	При нулевом среднем $\mathbf m=0$ ковариационная матрица совпадает с корреляционной: $\mathbf B=\mathbf K$. Поэтому плотность полностью выражается через корреляционную функцию:
	\[
	W_n(\mathbf x)=
	\frac{1}{(2\pi)^{n/2}\sqrt{\det\mathbf K}}
	\exp\left(-\frac12\mathbf x^T\mathbf K^{-1}\mathbf x\right).
	\]
	
	Основные свойства:
	
	Гауссовский процесс полностью определяется средним значением и ковариационной функцией. Все линейные преобразования гауссовского процесса и линейные комбинации гауссовских процессов сохраняют гауссовость. Условные плотности также имеют гауссову форму. Для гауссовских процессов некоррелированность отсчётов эквивалентна их статистической независимости.
	
	Для процесса с нулевым средним все нечётные моменты равны нулю, а чётные выражаются через ковариационную функцию суммированием по всем возможным попарным разбиениям (формула Иссерлиса--Вика):
	\[
	\boxed{
		\begin{aligned}
			\langle x_1x_2x_3x_4\rangle
			={}&B_x(t_1,t_2)B_x(t_3,t_4)\\
			&+B_x(t_1,t_3)B_x(t_2,t_4)\\
			&+B_x(t_1,t_4)B_x(t_2,t_3).
		\end{aligned}
	}
	\]
	В общем случае:
	\[
	\langle x(t_1)\cdots x(t_{2m})\rangle
	=
	\sum_{\text{всем попарным разбиениям}}
	\prod_{(i,k)}B_x(t_i,t_k).
	\]
	
	В частности, в одномерном случае $\langle x^{2m}(t)\rangle=(2m-1)!!\,\sigma_x^{2m}(t)$.
	
	\vspace{2mm}
	\textbf{Вопросы, не раскрытые в переданных материалах:} нет.
	
\end{document}